\documentclass[12pt,a4paper]{article}
\usepackage[utf8]{vietnam}
\usepackage{amsmath,amssymb}
\usepackage{tabularx}
\usepackage{tikz}
\usepackage{xcolor}
\usepackage[top=0.6cm, bottom=0.6cm, left=0.6cm, right=0.6cm]{geometry}
\usepackage[dethi]{ex_test}

\setlength{\parskip}{0pt}

% Dong ke cham mo chuan viet tay
\newcommand{\dongke}[1]{%
  \begingroup%
  \setlength{\parskip}{0pt}%
  \baselineskip=18.5pt%
  \vspace{1.0mm}%
  \foreach \i in {1,...,#1}{%
    \noindent{\color{gray!60}\dotfill}\par%
  }%
  \vspace{0.6mm}%
  \endgroup%
}

% Dinh nghia o bai tap tu luan trong bang cho Trang 1 va Trang 2
\newcommand{\oBaiTapMot}[3][4]{%
  \vspace{2.8mm}%
  \noindent\textbf{#2} #3\par%
  \dongke{#1}%
}
\newcommand{\oBaiTapHai}[3][5]{%
  \vspace{1.0mm}%
  \noindent\textbf{#2} #3\par%
  \dongke{#1}%
}

\begin{document}
\pagestyle{empty}
% =============================================================================
% ========================= MA DE 101 - TRANG 1/2 =============================
% =============================================================================
\noindent
\begin{tabularx}{\linewidth}{|X|c|}
\hline
\textbf{PHIẾU RÈN LUYỆN ĐẠI SỐ TOÁN 8 - TUẦN 11 - BUỔI 02} & \textbf{Mã đề: 101 - Trang 1/2} \\
\hline
\end{tabularx}

\vspace{2.2mm}

\noindent
\begin{tabularx}{\linewidth}{|p{3.2cm}|X|}
\hline
\centering\textbf{ĐIỂM SỐ} & \centering\textbf{LỜI NHẬN XÉT CỦA THẦY CÔ} \tabularnewline
\hline
\rule{0pt}{1.8cm} & \tabularnewline
\hline
\end{tabularx}

\vspace{2.2mm}
\noindent\textbf{Câu 1 (5,0 điểm). Thực hiện các phép tính sau:}
\vspace{1.0mm}

\noindent
\begin{tabularx}{\linewidth}{|X|X|}
\hline
\oBaiTapMot[4]{a) (0,4 điểm)}{$2x(3x + 4)$} & 
\oBaiTapMot[4]{b) (0,4 điểm)}{$3x(2x^2 - 4x + 1)$} \\
\hline
\oBaiTapMot[4]{c) (0,4 điểm)}{$(4x^2y - 3xy^2) + (2x^2y + 5xy^2)$} & 
\oBaiTapMot[4]{d) (0,4 điểm)}{$(5x^3 - 2x^2 + 4) - (2x^3 - 2x^2 - 5)$} \\
\hline
\oBaiTapMot[4]{e) (0,4 điểm)}{$(x + 3)(x + 5)$} & 
\oBaiTapMot[4]{f) (0,4 điểm)}{$(x - 2)(x + 6)$} \\
\hline
\oBaiTapMot[4]{g) (0,4 điểm)}{$(2x + 3)^2$} & 
\oBaiTapMot[4]{h) (0,4 điểm)}{$(3x - 2)^2$} \\
\hline
\oBaiTapMot[4]{i) (0,4 điểm)}{$(2x - 5)(2x + 5)$} & 
\oBaiTapMot[4]{j) (0,4 điểm)}{$(12x^4y^3 - 8x^3y^2 + 4x^2y) : 4x^2y$} \\
\hline
\oBaiTapMot[4]{k) (0,5 điểm)}{$(x + 4)^2 - x(x + 5)$} & 
\oBaiTapMot[4]{l) (0,5 điểm)}{$(2x^3 - 5x^2 + 6x - 2) : (2x - 1)$} \\
\hline
\end{tabularx}

\newpage
% =============================================================================
% ========================= MA DE 101 - TRANG 2/2 =============================
% =============================================================================
\noindent
\begin{tabularx}{\linewidth}{|X|c|}
\hline
\textbf{TOÁN 8 - TUẦN 11 - BUỔI 02 (RÈN LUYỆN ĐẠI SỐ)} & \textbf{Mã đề: 101 - Trang 2/2} \\
\hline
\end{tabularx}

\vspace{2.0mm}
\noindent\textbf{Câu 2 (5,0 điểm). Phân tích các đa thức sau thành nhân tử:}
\vspace{0.6mm}

\noindent
\begin{tabularx}{\linewidth}{|X|X|}
\hline
\oBaiTapHai[5]{a) (0,4 điểm)}{$3x^2 - 6x$} & 
\oBaiTapHai[5]{b) (0,4 điểm)}{$5x^2y - 10xy^2$} \\
\hline
\oBaiTapHai[5]{c) (0,4 điểm)}{$x(x - 3) + 5(x - 3)$} & 
\oBaiTapHai[5]{d) (0,4 điểm)}{$x^2 - 16$} \\
\hline
\oBaiTapHai[5]{e) (0,4 điểm)}{$4x^2 - 9$} & 
\oBaiTapHai[5]{f) (0,4 điểm)}{$x^2 - 6x + 9$} \\
\hline
\oBaiTapHai[5]{g) (0,4 điểm)}{$4x^2 - 12x + 9$} & 
\oBaiTapHai[5]{h) (0,4 điểm)}{$xy + 3y + 2x + 6$} \\
\hline
\oBaiTapHai[5]{i) (0,4 điểm)}{$xy - 2y + 3x - 6$} & 
\oBaiTapHai[5]{j) (0,4 điểm)}{$x^2 - 2xy + y^2 - 16$} \\
\hline
\oBaiTapHai[5]{k) (0,5 điểm)}{$x^2 - 5x + 6$} & 
\oBaiTapHai[5]{l) (0,5 điểm)}{$2x^3 - 8x$} \\
\hline
\end{tabularx}

\newpage
% =============================================================================
% ========================= MA DE 102 - TRANG 1/2 =============================
% =============================================================================
\noindent
\begin{tabularx}{\linewidth}{|X|c|}
\hline
\textbf{PHIẾU RÈN LUYỆN ĐẠI SỐ TOÁN 8 - TUẦN 11 - BUỔI 02} & \textbf{Mã đề: 102 - Trang 1/2} \\
\hline
\end{tabularx}

\vspace{2.2mm}

\noindent
\begin{tabularx}{\linewidth}{|p{3.2cm}|X|}
\hline
\centering\textbf{ĐIỂM SỐ} & \centering\textbf{LỜI NHẬN XÉT CỦA THẦY CÔ} \tabularnewline
\hline
\rule{0pt}{1.8cm} & \tabularnewline
\hline
\end{tabularx}

\vspace{2.2mm}
\noindent\textbf{Câu 1 (5,0 điểm). Thực hiện các phép tính sau:}
\vspace{1.0mm}

\noindent
\begin{tabularx}{\linewidth}{|X|X|}
\hline
\oBaiTapMot[4]{a) (0,4 điểm)}{$3x(2x + 5)$} & 
\oBaiTapMot[4]{b) (0,4 điểm)}{$2x(3x^2 - 5x + 2)$} \\
\hline
\oBaiTapMot[4]{c) (0,4 điểm)}{$(5x^2y - 4xy^2) + (x^2y + 6xy^2)$} & 
\oBaiTapMot[4]{d) (0,4 điểm)}{$(6x^3 - 3x^2 + 7) - (2x^3 - 3x^2 - 1)$} \\
\hline
\oBaiTapMot[4]{e) (0,4 điểm)}{$(x + 2)(x + 6)$} & 
\oBaiTapMot[4]{f) (0,4 điểm)}{$(x - 3)(x + 7)$} \\
\hline
\oBaiTapMot[4]{g) (0,4 điểm)}{$(3x + 2)^2$} & 
\oBaiTapMot[4]{h) (0,4 điểm)}{$(2x - 3)^2$} \\
\hline
\oBaiTapMot[4]{i) (0,4 điểm)}{$(3x - 4)(3x + 4)$} & 
\oBaiTapMot[4]{j) (0,4 điểm)}{$(15x^4y^3 - 9x^3y^2 + 3x^2y) : 3x^2y$} \\
\hline
\oBaiTapMot[4]{k) (0,5 điểm)}{$(x + 3)^2 - x(x + 4)$} & 
\oBaiTapMot[4]{l) (0,5 điểm)}{$(2x^3 - 7x^2 + 7x - 2) : (2x - 1)$} \\
\hline
\end{tabularx}

\newpage
% =============================================================================
% ========================= MA DE 102 - TRANG 2/2 =============================
% =============================================================================
\noindent
\begin{tabularx}{\linewidth}{|X|c|}
\hline
\textbf{TOÁN 8 - TUẦN 11 - BUỔI 02 (RÈN LUYỆN ĐẠI SỐ)} & \textbf{Mã đề: 102 - Trang 2/2} \\
\hline
\end{tabularx}

\vspace{2.0mm}
\noindent\textbf{Câu 2 (5,0 điểm). Phân tích các đa thức sau thành nhân tử:}
\vspace{0.6mm}

\noindent
\begin{tabularx}{\linewidth}{|X|X|}
\hline
\oBaiTapHai[5]{a) (0,4 điểm)}{$4x^2 - 8x$} & 
\oBaiTapHai[5]{b) (0,4 điểm)}{$6x^2y - 12xy^2$} \\
\hline
\oBaiTapHai[5]{c) (0,4 điểm)}{$x(x - 4) + 3(x - 4)$} & 
\oBaiTapHai[5]{d) (0,4 điểm)}{$x^2 - 25$} \\
\hline
\oBaiTapHai[5]{e) (0,4 điểm)}{$9x^2 - 16$} & 
\oBaiTapHai[5]{f) (0,4 điểm)}{$x^2 - 8x + 16$} \\
\hline
\oBaiTapHai[5]{g) (0,4 điểm)}{$9x^2 - 12x + 4$} & 
\oBaiTapHai[5]{h) (0,4 điểm)}{$xy + 4y + 3x + 12$} \\
\hline
\oBaiTapHai[5]{i) (0,4 điểm)}{$xy - 3y + 4x - 12$} & 
\oBaiTapHai[5]{j) (0,4 điểm)}{$x^2 - 2xy + y^2 - 25$} \\
\hline
\oBaiTapHai[5]{k) (0,5 điểm)}{$x^2 - 7x + 12$} & 
\oBaiTapHai[5]{l) (0,5 điểm)}{$3x^3 - 12x$} \\
\hline
\end{tabularx}

\newpage
% =============================================================================
% ========================= MA DE 103 - TRANG 1/2 =============================
% =============================================================================
\noindent
\begin{tabularx}{\linewidth}{|X|c|}
\hline
\textbf{PHIẾU RÈN LUYỆN ĐẠI SỐ TOÁN 8 - TUẦN 11 - BUỔI 02} & \textbf{Mã đề: 103 - Trang 1/2} \\
\hline
\end{tabularx}

\vspace{2.2mm}

\noindent
\begin{tabularx}{\linewidth}{|p{3.2cm}|X|}
\hline
\centering\textbf{ĐIỂM SỐ} & \centering\textbf{LỜI NHẬN XÉT CỦA THẦY CÔ} \tabularnewline
\hline
\rule{0pt}{1.8cm} & \tabularnewline
\hline
\end{tabularx}

\vspace{2.2mm}
\noindent\textbf{Câu 1 (5,0 điểm). Thực hiện các phép tính sau:}
\vspace{1.0mm}

\noindent
\begin{tabularx}{\linewidth}{|X|X|}
\hline
\oBaiTapMot[4]{a) (0,4 điểm)}{$4x(2x + 3)$} & 
\oBaiTapMot[4]{b) (0,4 điểm)}{$3x(3x^2 - 4x + 2)$} \\
\hline
\oBaiTapMot[4]{c) (0,4 điểm)}{$(3x^2y - 5xy^2) + (4x^2y + 7xy^2)$} & 
\oBaiTapMot[4]{d) (0,4 điểm)}{$(7x^3 - 4x^2 + 6) - (3x^3 - 4x^2 - 2)$} \\
\hline
\oBaiTapMot[4]{e) (0,4 điểm)}{$(x + 4)(x + 3)$} & 
\oBaiTapMot[4]{f) (0,4 điểm)}{$(x - 1)(x + 5)$} \\
\hline
\oBaiTapMot[4]{g) (0,4 điểm)}{$(2x + 5)^2$} & 
\oBaiTapMot[4]{h) (0,4 điểm)}{$(3x - 1)^2$} \\
\hline
\oBaiTapMot[4]{i) (0,4 điểm)}{$(4x - 3)(4x + 3)$} & 
\oBaiTapMot[4]{j) (0,4 điểm)}{$(16x^4y^3 - 12x^3y^2 + 4x^2y) : 4x^2y$} \\
\hline
\oBaiTapMot[4]{k) (0,5 điểm)}{$(x + 5)^2 - x(x + 7)$} & 
\oBaiTapMot[4]{l) (0,5 điểm)}{$(3x^3 - 7x^2 + 5x - 1) : (3x - 1)$} \\
\hline
\end{tabularx}

\newpage
% =============================================================================
% ========================= MA DE 103 - TRANG 2/2 =============================
% =============================================================================
\noindent
\begin{tabularx}{\linewidth}{|X|c|}
\hline
\textbf{TOÁN 8 - TUẦN 11 - BUỔI 02 (RÈN LUYỆN ĐẠI SỐ)} & \textbf{Mã đề: 103 - Trang 2/2} \\
\hline
\end{tabularx}

\vspace{2.0mm}
\noindent\textbf{Câu 2 (5,0 điểm). Phân tích các đa thức sau thành nhân tử:}
\vspace{0.6mm}

\noindent
\begin{tabularx}{\linewidth}{|X|X|}
\hline
\oBaiTapHai[5]{a) (0,4 điểm)}{$5x^2 - 10x$} & 
\oBaiTapHai[5]{b) (0,4 điểm)}{$4x^2y - 12xy^2$} \\
\hline
\oBaiTapHai[5]{c) (0,4 điểm)}{$x(x - 2) + 6(x - 2)$} & 
\oBaiTapHai[5]{d) (0,4 điểm)}{$x^2 - 36$} \\
\hline
\oBaiTapHai[5]{e) (0,4 điểm)}{$16x^2 - 25$} & 
\oBaiTapHai[5]{f) (0,4 điểm)}{$x^2 - 10x + 25$} \\
\hline
\oBaiTapHai[5]{g) (0,4 điểm)}{$4x^2 + 20x + 25$} & 
\oBaiTapHai[5]{h) (0,4 điểm)}{$xy + 5y + 2x + 10$} \\
\hline
\oBaiTapHai[5]{i) (0,4 điểm)}{$xy - 4y + 2x - 8$} & 
\oBaiTapHai[5]{j) (0,4 điểm)}{$x^2 - 2xy + y^2 - 36$} \\
\hline
\oBaiTapHai[5]{k) (0,5 điểm)}{$x^2 - 6x + 8$} & 
\oBaiTapHai[5]{l) (0,5 điểm)}{$2x^3 - 18x$} \\
\hline
\end{tabularx}

\newpage
% =============================================================================
% ========================= MA DE 104 - TRANG 1/2 =============================
% =============================================================================
\noindent
\begin{tabularx}{\linewidth}{|X|c|}
\hline
\textbf{PHIẾU RÈN LUYỆN ĐẠI SỐ TOÁN 8 - TUẦN 11 - BUỔI 02} & \textbf{Mã đề: 104 - Trang 1/2} \\
\hline
\end{tabularx}

\vspace{2.2mm}

\noindent
\begin{tabularx}{\linewidth}{|p{3.2cm}|X|}
\hline
\centering\textbf{ĐIỂM SỐ} & \centering\textbf{LỜI NHẬN XÉT CỦA THẦY CÔ} \tabularnewline
\hline
\rule{0pt}{1.8cm} & \tabularnewline
\hline
\end{tabularx}

\vspace{2.2mm}
\noindent\textbf{Câu 1 (5,0 điểm). Thực hiện các phép tính sau:}
\vspace{1.0mm}

\noindent
\begin{tabularx}{\linewidth}{|X|X|}
\hline
\oBaiTapMot[4]{a) (0,4 điểm)}{$5x(2x + 3)$} & 
\oBaiTapMot[4]{b) (0,4 điểm)}{$4x(2x^2 - 3x + 1)$} \\
\hline
\oBaiTapMot[4]{c) (0,4 điểm)}{$(6x^2y - 2xy^2) + (2x^2y + 5xy^2)$} & 
\oBaiTapMot[4]{d) (0,4 điểm)}{$(5x^3 - 5x^2 + 8) - (x^3 - 5x^2 - 2)$} \\
\hline
\oBaiTapMot[4]{e) (0,4 điểm)}{$(x + 1)(x + 7)$} & 
\oBaiTapMot[4]{f) (0,4 điểm)}{$(x - 4)(x + 8)$} \\
\hline
\oBaiTapMot[4]{g) (0,4 điểm)}{$(3x + 1)^2$} & 
\oBaiTapMot[4]{h) (0,4 điểm)}{$(2x - 5)^2$} \\
\hline
\oBaiTapMot[4]{i) (0,4 điểm)}{$(5x - 2)(5x + 2)$} & 
\oBaiTapMot[4]{j) (0,4 điểm)}{$(18x^4y^3 - 12x^3y^2 + 6x^2y) : 6x^2y$} \\
\hline
\oBaiTapMot[4]{k) (0,5 điểm)}{$(x + 4)^2 - x(x + 6)$} & 
\oBaiTapMot[4]{l) (0,5 điểm)}{$(2x^3 - 9x^2 + 13x - 6) : (2x - 3)$} \\
\hline
\end{tabularx}

\newpage
% =============================================================================
% ========================= MA DE 104 - TRANG 2/2 =============================
% =============================================================================
\noindent
\begin{tabularx}{\linewidth}{|X|c|}
\hline
\textbf{TOÁN 8 - TUẦN 11 - BUỔI 02 (RÈN LUYỆN ĐẠI SỐ)} & \textbf{Mã đề: 104 - Trang 2/2} \\
\hline
\end{tabularx}

\vspace{2.0mm}
\noindent\textbf{Câu 2 (5,0 điểm). Phân tích các đa thức sau thành nhân tử:}
\vspace{0.6mm}

\noindent
\begin{tabularx}{\linewidth}{|X|X|}
\hline
\oBaiTapHai[5]{a) (0,4 điểm)}{$2x^2 - 6x$} & 
\oBaiTapHai[5]{b) (0,4 điểm)}{$7x^2y - 14xy^2$} \\
\hline
\oBaiTapHai[5]{c) (0,4 điểm)}{$x(x - 5) + 2(x - 5)$} & 
\oBaiTapHai[5]{d) (0,4 điểm)}{$x^2 - 49$} \\
\hline
\oBaiTapHai[5]{e) (0,4 điểm)}{$25x^2 - 4$} & 
\oBaiTapHai[5]{f) (0,4 điểm)}{$x^2 + 6x + 9$} \\
\hline
\oBaiTapHai[5]{g) (0,4 điểm)}{$9x^2 + 6x + 1$} & 
\oBaiTapHai[5]{h) (0,4 điểm)}{$xy + 2y + 4x + 8$} \\
\hline
\oBaiTapHai[5]{i) (0,4 điểm)}{$xy - 5y + 3x - 15$} & 
\oBaiTapHai[5]{j) (0,4 điểm)}{$x^2 - 2xy + y^2 - 49$} \\
\hline
\oBaiTapHai[5]{k) (0,5 điểm)}{$x^2 - 8x + 15$} & 
\oBaiTapHai[5]{l) (0,5 điểm)}{$5x^3 - 20x$} \\
\hline
\end{tabularx}
\end{document}
